\documentclass[11pt]{beamer}
\usepackage[utf8]{inputenc}
\usepackage[T1]{fontenc}
\usepackage{lmodern}
\usepackage[]{babel}
\usetheme{madrid}
\usepackage{fontspec}
\usepackage{color}
\usepackage{xcolor}
\setsansfont{Times New Roman}
\usepackage{graphics}
\usepackage{graphicx}
\usepackage{float}
\usepackage{subfigure}
\usepackage{multicol}
\usepackage{fancyhdr}
\usepackage{lipsum}
\usepackage{tikz}
\usepackage{wallpaper}
\usepackage{wrapfig}
\usepackage{multicol}
\begin{document}
	\author[GG BOY]{GOOD GUYS\inst{1} }%作者
	\title{Federated Learning}%大标题
	\subtitle{Distributed \& Encryption}%小标题
	\logo{\includegraphics[scale=0.08]{777}}%logo，默认每页都有，处在右下角
	\institute[GDUT]{\inst{1} \textit{School of Mathematics and Statistics,\\ Guangdong University of Technology} }%单位
	%\date{}
	%\subject{}
	%\setbeamercovered{transparent}
	\setbeamertemplate{navigation symbols}{}%原本Madrid主题每页下面都有一些信息，这个命令是去除下方的信息，不过我用数院长的logo挡住了，删不删也看不到。
	\setbeamertemplate{itemize }{circle}%定义item的形状是小圆圈
	\newcommand\Background{
            \begin{tikzpicture}[remember picture,overlay]
			\node[inner sep=0pt,outer sep=0pt,opacity=1]
			at (5.7,-4)
			{\includegraphics[scale=0.23]{999}};
		\end{tikzpicture}
        }%定义了背景，想在哪页加背景，只需在哪页加一句\Background即可
	%\setbeamertemplate{background canvas}{\includegraphics[scale=0.19]{999}}%这句命令也是定义了背景，但是没法设置指定页，必须全部页都有背景，要是不嫌背景花直接用这句命令添加背景省事。
	\setbeamercolor{frametitle}{bg=gray}%定义frametitle的背景是灰色，也就是每一页上面那栏
	\addtobeamertemplate{footline}{}{
        \begin{tikzpicture}
			{\-0.8cm\includegraphics[scale=0.2]{000}}%设置了每页footline，也就是最下面加数院的长logo
	\end{tikzpicture}}
\begin{frame}[plain]
	\Background%要先生成背景再生成title，要不然背景会覆盖在title页面上
	\maketitle
\end{frame}

\begin{frame}
	\Background
	\tableofcontents
\end{frame}

\begin{frame}
	\frametitle{\textsc{Idealized}}
	\section{Background}
	\subsection{Idealized}
	\subsection{Realistic}
		aaaaaaaa
	\begin{examples}
		bbbbbbbbbb
	\end{examples}
		
	\begin{figure}[t]
		\centering
		\subfigure{\includegraphics[height=3cm ]{111}}
		\subfigure{\includegraphics[width=3cm]{222}}
	\end{figure}
	bbbbbbbbbb
\end{frame}
	
\begin{frame}
	\frametitle{\textsc{Idealized}}
	aaaaaaaaaaa
	\begin{block}{Question}
		bbbbbbb
	\end{block}
	NO NO NO
	\begin{alertblock}{Important}
		aaaaaa
	\end{alertblock}
		 
\end{frame}

\begin{frame}
	\section{Introduction}
	\subsection{Concept}
	\frametitle{\textsc{Concept}}
\end{frame}

\begin{frame}
	\section{Classification}
	\subsection{Horizontal federated learning}
	\frametitle{\textsc{Horizontal Federated Learning}}
	\begin{itemize}
		\item a
		\item b
		\item Eac
		\item d
		\item e
	\end{itemize}
\end{frame}

\begin{frame}
	\subsection{Longitudinal federated learning}
	\frametitle{\textsc{Longitudinal Federated Learning}}
	\begin{block}{Definition}
		aaa
	\end{block}
	\begin{example}
		ddd
	\end{example}
\end{frame}

\begin{frame}
	\frametitle{\textsc{Longitudinal Federated Learning}}
	
	
\end{frame}

\begin{frame}
	\subsection{Transfer federated learning}
	\frametitle{\textsc{Transfer Federated Learning}}
\end{frame}

\begin{frame}
	\frametitle{\textsc{References}}
	\Background
\end{frame}
\begin{frame}[plain]
	\Background
	\begin{center}
		{\Huge $\mathsf{THANK YOU}$}\\
		\vspace{5pt}
		DA SHA BI $^{1}$ \\
		\vspace{5pt}
		$^{1}$ \textit{School of Mathematics and Statistics,\\ Guangdong University of Technology}
	\end{center}
	
\end{frame}
\end{document}